% Punkte und Strecken im Koordinatensystem · Körper und Drehungen · Prüfungs-Fokus Abitur GK – bank/punkte-und-strecken-im-koordinatensystem/e5.jsonl (zusammenbau v0.6, Prüfungs-Fokus)
\einheitenkopf[e5]{Einheit 5 · Körper im Koordinatensystem und Drehungen}
\zweigzeile{Abi GK ×7}
\setcounter{aufgabe}{0}

% Anlauf (ohne Prüfkennung)
\begin{aufgabe}{Körper und Drehungen – Anlauf}
\begin{teile}
\teil Ein Quader steht auf der $x_1x_2$-Ebene; seine Deckenecke $G$ hat $x_3 = 4$. Was sagt das über die Deckfläche? \\ \kreuz{Sie liegt in der $x_1x_2$-Ebene.} \\ \kreuz{Sie liegt parallel zur $x_1x_2$-Ebene in der Höhe $4$.} \\ \kreuz{Sie steht senkrecht auf der $x_1x_2$-Ebene.}
\end{teile}
\end{aufgabe}

\begin{aufgabe}{Körper und Drehungen – Prüfungsaufgaben}
\begin{teile}
\teil Gerades Prisma $ABCDEF$ ($D$ über $A$, $E$ über $B$, $F$ über $C$) mit $A(2 | 1 | 0)$, $B(7 | 1 | 0)$, $C(2 | 4 | 0)$, $D(2 | 1 | 3)$: $F$? \hfill (A26) \\ F( \leerfeld | \leerfeld | \leerfeld )
\teil Quader $ABCDEFGH$ (Boden $ABCD$, $E$ über $A$, $F$ über $B$, $G$ über $C$, $H$ über $D$) mit $A(2 | 1 | 0)$, $B(6 | 1 | 0)$, $E(2 | 1 | 5)$, $H(2 | 4 | 5)$: $G$? \hfill (A24) \\ G( \leerfeld | \leerfeld | \leerfeld )
\teil Quader $ABCDEFGH$ (Boden $ABCD$, $E$ über $A$, $F$ über $B$, $G$ über $C$, $H$ über $D$) mit $A(0 | 2 | 1)$, $B(4 | 2 | 1)$, $E(0 | 2 | 4)$, $H(0 | 7 | 4)$: $G$? \hfill (A24) \\ G( \leerfeld | \leerfeld | \leerfeld )
\teil Gerades Prisma $ABCDEF$ ($D$ über $A$, $E$ über $B$, $F$ über $C$) mit $A(4 | 1 | 2)$, $B(4 | 6 | 2)$, $C(8 | 1 | 2)$, $D(4 | 1 | 7)$: $F$? \hfill (A26) \\ F( \leerfeld | \leerfeld | \leerfeld )
\teil Eine Pyramide hat als Grundfläche ein rechtwinkliges Dreieck mit den Katheten $6$ cm und $8$ cm; die Spitze liegt $5$ cm senkrecht über der Ecke mit dem rechten Winkel. Wähle geeignete Koordinaten der vier Ecken (1 LE = 1 cm). \hfill (A20)
\teil Eine Halle ist ein Quader mit quadratischem Grundriss (Seite $15$ m) und $7$ m Höhe. $A$ liegt im Ursprung, $B$ auf der $x_1$-Achse, $D$ auf der $x_2$-Achse, die Deckenecken $E$, $F$, $G$, $H$ über $A$, $B$, $C$, $D$ (1 LE = 1 m). Koordinaten von $G$ und $H$? \hfill (A20)
\teil Ein Würfel mit Kante $2$ hat eine Ecke im Ursprung, alle Ecken haben Koordinaten $\ge 0$. Er wird um $(-1 | -1 | 1)$ verschoben. Welche Ecke hat danach $x_1 < 0$, $x_2 > 0$ und $x_3 > 2$? \hfill (A24) \\ ( \leerfeld | \leerfeld | \leerfeld )
\teil Ein Quader mit den Kanten $4$, $2$ und $3$ entlang der $x_1$-, $x_2$- und $x_3$-Achse hat eine Ecke im Ursprung, alle Ecken haben Koordinaten $\ge 0$. Er wird um $(-2 | 1 | -1)$ verschoben. Welche Ecke hat danach $x_1 > 0$, $x_2 < 2$ und $x_3 < 0$? \hfill (A24) \\ ( \leerfeld | \leerfeld | \leerfeld )
\teil Ein Quader mit den Kanten $3$, $5$ und $2$ entlang der $x_1$-, $x_2$- und $x_3$-Achse hat eine Ecke im Ursprung, alle Ecken haben Koordinaten $\ge 0$. Er wird um $(-2 | -3 | -1)$ verschoben. Welche Ecke hat danach $x_1 > 0$, $x_2 < 0$ und $x_3 > 0$? \hfill (A24) \\ ( \leerfeld | \leerfeld | \leerfeld )
\teil Eine Pyramide hat die Grundfläche $A(4 | -1 | 2)$, $B(1 | 3 | 2)$, $C(-1 | 0 | 2)$ und die Spitze $S(2 | 1 | 7)$. Begründe: Die Grundfläche ist parallel zur $x_1x_2$-Ebene. Höhe der Pyramide? \hfill (A18) \\ h = \leerfeld
\teil Eine Pyramide hat die Grundfläche $A(2 | 2 | -1)$, $B(5 | 0 | -1)$, $C(3 | 5 | -1)$ und die Spitze $S(3 | 2 | 3)$. Begründe: Die Grundfläche ist parallel zur $x_1x_2$-Ebene. Höhe der Pyramide? \hfill (A18) \\ h = \leerfeld
\teil Eine Pyramide hat die Grundfläche $A(1 | 0 | 3)$, $B(4 | 2 | 3)$, $C(0 | 5 | 3)$ und die Spitze $S(2 | 2 | 9)$. Begründe: Die Grundfläche ist parallel zur $x_1x_2$-Ebene. Höhe der Pyramide? \hfill (A18) \\ h = \leerfeld
\end{teile}
\end{aufgabe}

\begin{aufgabe}{Körper und Drehungen – Prüfungsaufgaben – weiter}
\begin{teile}
\teil Ein Dachgeschoss hat einen rechteckigen Boden von $10$ m mal $8$ m. Die Dachflächen steigen von den beiden Traufwänden (Höhe $0{,}5$ m) gleichmäßig zum First über der Mitte (Höhe $5{,}5$ m). Flächen mit weniger als $2$ m Raumhöhe zählen nicht zur Wohnfläche. Welcher Anteil des Bodens zählt nicht? \hfill (A19) \\ \leerfeld \%
\teil Ein Dachgeschoss hat einen rechteckigen Boden von $14$ m mal $10$ m. Die Dachflächen steigen von den beiden Traufwänden (Höhe $1$ m) gleichmäßig zum First über der Mitte (Höhe $5$ m). Flächen mit weniger als $1{,}8$ m Raumhöhe zählen nicht zur Wohnfläche. Welcher Anteil des Bodens zählt nicht? \hfill (A19) \\ \leerfeld \%
\teil Ein Dachgeschoss hat einen rechteckigen Boden von $15$ m mal $10$ m. Die Dachflächen steigen von den beiden Traufwänden (Höhe $0{,}4$ m) gleichmäßig zum First über der Mitte (Höhe $2{,}9$ m). Flächen mit weniger als $1$ m Raumhöhe zählen nicht zur Wohnfläche. Welcher Anteil des Bodens zählt nicht? \hfill (A19) \\ \leerfeld \%
\teil Das Dreieck $A(5 | 3 | 0)$, $B(3 | 5 | 0)$, $C(3 | 3 | 4)$ wird um die Kante $AB$ gedreht, bis $C$ in der $x_1x_2$-Ebene liegt, mit $x_1 < 4$. Koordinaten des gedrehten Punktes $C'$? \hfill (A23) \\ C'( \leerfeld | \leerfeld | \leerfeld )
\teil Das Dreieck $A(2 | 1 | 0)$, $B(2 | 9 | 0)$, $C(5 | 5 | 4)$ wird um die Kante $AB$ gedreht, bis $C$ in der $x_1x_2$-Ebene liegt, mit $x_1 > 2$. Koordinaten des gedrehten Punktes $C'$? \hfill (A23) \\ C'( \leerfeld | \leerfeld | \leerfeld )
\teil Eine Pyramide hat den quadratischen Boden $ABCD$ mit der Seite $4$ in der $x_1x_2$-Ebene, $A$ im Ursprung, $B$ auf der $x_1$-Achse, $D$ auf der $x_2$-Achse; die Spitze $S$ liegt $6$ senkrecht über $A$. Ein Würfel hat drei Seitenflächen in den Koordinatenebenen und eine Ecke auf der Kante $CS$. Kantenlänge des Würfels? Begründe, dass kein Punkt des Würfels außerhalb der Pyramide liegt. \hfill (A24)
\teil Eine Pyramide hat den quadratischen Boden $ABCD$ mit der Seite $3$ in der $x_1x_2$-Ebene, $A$ im Ursprung, $B$ auf der $x_1$-Achse, $D$ auf der $x_2$-Achse; die Spitze $S$ liegt $6$ senkrecht über $A$. Ein Würfel hat drei Seitenflächen in den Koordinatenebenen und eine Ecke auf der Kante $CS$. Kantenlänge des Würfels? Begründe, dass kein Punkt des Würfels außerhalb der Pyramide liegt. \hfill (A24)
\teil Ein Körper hat den rechteckigen Boden $A(2 | 1 | 0)$, $B(6 | 1 | 0)$, $C(6 | 7 | 0)$, $D(2 | 7 | 0)$ und die Firstkante $P(4 | 2 | 3)$, $Q(4 | 5 | 3)$; die Dreiecke $ABP$ und $CDQ$ und die Vierecke $BCQP$ und $DAPQ$ begrenzen ihn. Die Ebenen $x_2 = s$ mit $1 < s < 7$ schneiden ihn. Wie viele Ecken hat die Schnittfläche? Für welche $s$ haben die Schnittflächen dieselbe Form und denselben Flächeninhalt? \hfill (A24)
\teil Ein Körper hat den rechteckigen Boden $A(0 | 1 | 0)$, $B(6 | 1 | 0)$, $C(6 | 9 | 0)$, $D(0 | 9 | 0)$ und die Firstkante $P(3 | 3 | 4)$, $Q(3 | 6 | 4)$; die Dreiecke $ABP$ und $CDQ$ und die Vierecke $BCQP$ und $DAPQ$ begrenzen ihn. Die Ebenen $x_2 = s$ mit $1 < s < 9$ schneiden ihn. Wie viele Ecken hat die Schnittfläche? Für welche $s$ haben die Schnittflächen dieselbe Form und denselben Flächeninhalt? \hfill (A24)
\teil Ein Würfel $ABCDEFGH$ mit Kante $4$ hat die Ecke $A$ im Ursprung, $B$ auf der $x_1$-Achse, $D$ auf der $x_2$-Achse; $E$ liegt über $A$, $G$ über $C$. Die Ebene durch $E$, $G$ und $S_t(t | 0 | 0)$ mit $0 \le t \le 4$ schneidet ihn. Eckenzahl der Schnittfläche in Abhängigkeit von $t$? Für welche $t$ hat sie genau zwei Symmetrieachsen? \hfill (A24)
\teil Ein Quader $ABCDEFGH$ mit quadratischem Boden (Seite $6$) und Höhe $3$ hat die Ecke $A$ im Ursprung, $B$ auf der $x_1$-Achse, $D$ auf der $x_2$-Achse; $E$ liegt über $A$, $G$ über $C$. Die Ebene durch $E$, $G$ und $S_t(t | 0 | 0)$ mit $0 \le t \le 6$ schneidet ihn. Eckenzahl der Schnittfläche in Abhängigkeit von $t$? Für welche $t$ hat sie genau zwei Symmetrieachsen? \hfill (A24)
\end{teile}
\end{aufgabe}

\medskip
\uebersichtskasten{Ecken ablesen: beim Quader und beim geraden Prisma liegt jede Deckenecke senkrecht über ihrer Grundecke – gleiche erste und zweite Koordinate, die dritte aus der Höhe; allgemein ist jede Ecke Nachbarecke plus Kantenvektor. \\ Prisma mit A(0 | 0 | 0), B(6 | 0 | 0), C(0 | 4 | 0), D(0 | 0 | 3): F liegt senkrecht über C in der Höhe von D, also F(0 | 4 | 3). \\ Koordinaten wählen: rechte Winkel auf die Achsen legen, eine Ecke in den Ursprung – ein nur durch Längen und rechte Winkel beschriebener Körper bekommt so die einfachsten Koordinaten. \\ Höhe und Symmetrie: liegt die Grundfläche parallel zu einer Koordinatenebene (gleiche dritte Koordinate aller Grundecken), ist die Höhe eine Koordinatendifferenz; ist der Körper symmetrisch zu einer Koordinatenebene, wechselt beim Partnerpunkt genau eine Koordinate das Vorzeichen. \\ Drehung: der gedrehte Punkt läuft auf einem Kreis um seinen Lotfußpunkt auf der Drehachse; der Radius ist der Abstand zur Achse, und in der Zielebene trägt man denselben Abstand senkrecht zur Achse ab – das Vorzeichen liefert die Bedingung der Aufgabe.}
